\documentclass[10pt]{article}

\usepackage[utf8]{inputenc}

\usepackage[T1]{fontenc}

\usepackage{amsmath}

\usepackage{amsfonts}

\usepackage{amssymb}

\usepackage[version=4]{mhchem}

\usepackage{stmaryrd}

\usepackage{bbold}

\usepackage{graphicx}

\usepackage[export]{adjustbox}

\graphicspath{ {./images/} }



\begin{document}

новение автоколебаний). Впервые достаточные условия возникновения автоколебаний в системе (*) доказал А. Льенар (A. Liénard, 1928).



Л. у. тесно связано с Рэлея уравнением. Важным частным случаем Л. у. является Ван дер Поля уравнение. Вместо уравнения (*) часто удобно рассматривать систему



\[

x^{\prime}=v, v^{\prime}=-x-f(x) v

\]



(автоколебательному процессу в системе (*) адекватен устойчивый предельный цикл на фазовой плоскости $x, v)$ или эквивалентное ей уравнение



\[

\frac{d v}{d x}=\frac{-x-f(x) v}{v} .

\]



Если ввести новую переменную



\[

y=x^{\prime}+F(x)

\]



где



\[

F(x)=\int_{0}^{x} f(\xi) d \xi

\]



то уравнение (*) переходит в систему



\[

x^{\prime}=y-F(x), y^{\prime}=-x .

\]



Более общими, чем Л. у., являются уравнения



\[

\begin{gathered}

x^{\prime \prime}+f(x) x^{\prime}+g(x)=0, \\

x^{\prime \prime}+\varphi\left(x, x^{\prime}\right) x^{\prime}+g(x)=0 .

\end{gathered}

\]



Основной интерес представляет выяснение возможно более широких достаточных условий, при к-рых эти уравнения имеют единственное устойчивое периодич. решение. Подробно изучалось также неоднородное Л. у.



и его обобщения.



\[

x^{\prime \prime}+f(x) x^{\prime}+x=e(t)

\]



Лит.: [1] Андрон о в А. А., В итт А. А., Хайкин С. Э., Теория колебаний, 2 изд., М., 1959; [2] Сансоне Дж., Обыкновенные дифференциальные уравнения, пер. с итал., т. 2, М., 1954; [3] Л е фше ц С., Геометрическая теория дифференциальных уравнений, пер. с англ., М., 1961; [4] Рейсс и г Р., Сансоне Г., Конти Р., Качественная теория нелинейных дифференциальных уравнений, пер. с нем., М., 1974. Н. Х. Розов. ЛЮДЕРСА - ПАУЛИ ТЕОРЕМА - см. СРТ-теореМа. ЛюКСОH - см. Тахион.



ЛЮСТЕРНИКА - ШНИРЕЛЬМАНА КАТЕГОРИЯ см. Вариационное исчисление в целом.



ЛЯПУНОВА ЛЕММА - см. Ляпунова функция.



ЛЯПУНОВА ПОКАЗАТЕЛЬ - сМ. Хаос.



ЛЯПУНОВА TEOPEMA об устойчивости - см. $\boldsymbol{y}_{c}$ тойчивость по Ляпунову.



ЛЯПУНОВА ФУНКЦИЯ - функция, определяемая следующим образом. Пусть $x_{0}-$ не под в ижная точка системы дифференциальных уравнений



\[

\dot{x}=f(x, t)

\]



\begin{itemize}

  \item (то есть $f\left(x_{0}, t\right) \equiv 0$ ), где отображение $f(x, t): U \times \mathbb{R}^{+} \rightarrow \mathbb{R}^{n}$ непрерывно и непрерывно дифференцируемо по $x$ (здесь $U$ - нек-рая окрестность точки $x_{0}$ в $\mathbb{R}^{n}$ ); в координатах эта система записывается в виде

\end{itemize}



\[

\dot{x}^{i}=f^{i}\left(x^{1}, \ldots, x^{n}, t\right), i=1, \ldots, n .

\]



Функц и ей Л я пун о в а называется дифференцируемая функция $V(x): U \rightarrow \mathbb{R}$, обладающая свойствами:



\begin{enumerate}

  \item $V(x)>0$ при $x \neq x_{0}$;



  \item $V\left(x_{0}\right)=0$



  \item $0 \geqslant \frac{d V(x)}{d x} f(x, t)=\sum_{i=1}^{n} \frac{\partial V\left(x^{1}, \ldots, x^{n}\right)}{\partial x^{i}} f^{i}\left(x^{1}, \ldots, x^{n}, t\right)$.



\end{enumerate}



Функция $V(x)$ введена А. М. Ляпуновым (см. [1]).



Имеет место ле м м а Л я п у н о ва: если Л. ф. существует, то неподвижная точка устойчива по Ляпунову. На этой лемме основан один из методов исследования устойчивости.



Jum.: [1] Л я п у н о в А. М., Собр. соч., т. 2, М.-Л., 1956, с. 7-263; [2] Б арба ш и н Е. А., Функции Ляпунова, М., 1970.



В. М. Миллионциков.



ЛЯПУНОВА - ПЕРРОНА ТЕОРЕМА - см. Центральное многообразие.



\begin{center}

\includegraphics[max width=\textwidth]{2023_07_08_ad66712331f78cc2ab25g-1}

\end{center}



MATHИTHOE УПOPЯДОЧЕНИЕ - см. Дальний и ближний порядок.



МАГНИТНОЙ ГИДРОДИНАМИКИ УРАВНЕНИЯ - УРавнения движения электропроводной жидкости или ионизованного газа (плазмы) в магнитном поле. При этом индуцируется электрич. поле и возникает электрич. ток. В свою очередь, взаимодействие тока с магнитным полем вызывает изменение движения жидкости и изменение магнитного поля. Магнитная гидродинамика в широком смысле - это раздел механики сплошных сред, изучающий движение электропроводных газов и жидкостей при наличии магнитного поля. При этом состояние среды характеризуется функциями распределения по скоростям $f_{\alpha}(\boldsymbol{v}, \boldsymbol{r}, t)$ для каждого сорта частиц $\alpha$, а также значениями электрич. и магнитного полей во всех точках пространства и в каждый момент времени. Если характерная для рассматриваемых процессов частота $\omega$ значительно меньше частоты столкновений отдельных частиц $v, \omega \ll v$, то каждый сорт частиц описывают гидродинамически. При этом состояние частиц сорта $\alpha$ описывается плотностью



\[

n_{\alpha}(r, t)=\int f_{\alpha}(v, r, t) d v

\]



и средней скоростью



\[

v_{\alpha}(r, t)=\frac{1}{n_{\alpha}} \int v f_{\alpha}(v, r, t) d v .

\]



Это - м ного ж ид к о с т а я магнитная гидродинамика. В частности, полностью ионизованной плазме, состоящей из электронов и одного сорта ионов, соответствует д в у $\mathrm{x}$ ж и д к о с т н а я магнитная гидродинамика.



Если характерная частота $\omega$ значительно меньше частоты обмена энергией между электроном и ионом $v_{e i}$, то давления электронов $p_{e}$ и ионов $p_{i}$ успевают выровняться и плазма характеризуется единым давлением $p$. Из-за того, что электрон значительно легче иона, плотность массы $\rho$ и скорость среды $v$ определяются ионами, а плотность тока $\boldsymbol{j}$ - электронами. Такую плазму, рассматриваемую как единое целое, называют п р о с т о й. Однако уравнения движения простой плазмы, дополненные уравнениями Максвелла, все еще слишком сложны. Решающий шаг в их упрощении был сделан Г. Альфвеном (H. Alfven,-1942). Он показал,





\end{document}